\begin{claim}{C.46}{Proved}{Two boundary foliations admit a compatible intrinsic contraction}
Use the actual based null-homotopic plane loop $f_t$, its generated group $G$, the
open orbit $O$, and $L,x,J$ of C.42. For every
$z\in\operatorname{Fix}(J)\cap\partial O$ with nontrivial
intrinsic line or circle orbit, lift its quotient $\Phi_z$
to the common source universal covering $p:X\to O$ as in
C.42. Let $\mathscr F$ be the set of fiber partitions of
these real lifts; changing their target coordinates or
their lift constants does not change this set.
For $|\mathscr F|=2$, the loop $f_t|_O$ has a based
compact-open null-homotopy
\[
 K_{s,t}\in\operatorname{Homeo}(O),\qquad
 K_{0,t}=f_t|_O,\quad K_{1,t}=K_{s,0}=K_{s,1}=\id.
                                                        \tag{C.46.1}
\]
It commutes with every plane homeomorphism commuting
with all $f_t$ and preserving $O$. For every indicated
point $z$, it also satisfies the exact quotient relation
\[
                 \Phi_z K_{s,t}=Q^{B_z}_{s,t}\Phi_z,
                  \qquad B_z=Gz,                       \tag{C.46.2}
\]
where $Q^{B_z}$ is the natural median contraction of the
original restricted loop on that intrinsic line or
circle, as in C.27--C.29. Thus all the specified
boundary quotient relations hold simultaneously.
The asserted homotopy is intrinsic to $O$; continuity
at its ambient frontier is a further assertion.
\end{claim}

\begin{proof}
Choose representatives $F_1,F_2:X\to\R$ of the two
partitions. C.45 supplies a homeomorphism from $X$
onto one of its three domains $D$, with these two
coordinates changed only by independent line
homeomorphisms and a possible exchange. Lift the
isotopy $f_t|_O$ from the identity to $a_t$ on $X$.
For each original boundary target, lift $f_t|_{B_i}$
to its based real path $b_{i,t}$. Covering uniqueness
gives
\[
                       F_i a_t=b_{i,t}F_i.              \tag{C.46.3}
\]
Each $b_{i,t}$ is increasing. For a line target its
endpoint is the identity. For a circle target C.29
identifies its endpoint translation with the winding
of the original ambient-null plane loop and proves
that translation zero. Hence both $b_{i,1}$ are the
identity. Since the pair $(F_1,F_2)$ is injective,
$a_1=\id$. We have obtained an actual based loop in
the corresponding coordinate-pair group (C.45.3),
without discarding a source or target deck shift.

In these coordinates write $a_t=(g_{1,t},g_{2,t})$.
Apply C.27 to each coordinate loop and set
\[
\begin{split}
 m_{i,r}(u)&=\min_{0\leq v\leq r}g_{i,v}(u),\qquad
 M_{i,r}(u)=\max_{0\leq v\leq r}g_{i,v}(u),\\
 R_{i,s,t}(u)&=\operatorname{med}
       \bigl(m_{i,1-s}(u),g_{i,t}(u),M_{i,1-s}(u)\bigr),\\
 \widetilde K_{s,t}(u,v)&=(R_{1,s,t}(u),R_{2,s,t}(v)).
\end{split}                                             \tag{C.46.4}
\]
The coordinate maps and their inverses are jointly
continuous by C.27 and C.31. They preserve the actual
domain $D$. For $D=\R^2$ this follows directly from
their being line homeomorphisms. For $D=\{u<v\}$,
C.45 gives $g_{1,t}=g_{2,t}$, so the two new coordinate
maps are identical and preserve the strict inequality.
For $D=\{u<v<u+1\}$ the same equality holds and each
$g_{i,t}$ commutes with translation by one. Minima,
maxima and the median preserve that equation, so the
new equal coordinate maps preserve both inequalities.
The inverse family has the same property. Formula
(C.46.4) is therefore a based null-homotopy in
$\operatorname{Homeo}(X)$ with the asserted endpoint
identities.

We establish the operator comparison before descending.
Let $h$ commute with all $f_t$ and preserve $O$.
It commutes with $L$ by continuity of composition and
inversion. The stabilizers of $x$ and $hx$ in $L$
are equal: an element fixes one precisely when it
fixes the other, by commutation with $h$ and $h^{-1}$.
Choose $a_h\in L$ with $a_hx=hx$. The exact stabilizer
identity gives $a_hJa_h^{-1}=J$. For every specified
$z$ and every $l\in L$, evaluation gives
\[
 \Phi_z h(lx)=\Phi_z(l a_hx)=l a_hz
                         =\Phi_{a_hz}(lx).             \tag{C.46.5}
\]
The point $a_hz$ lies in $\operatorname{Fix}(J)\cap
\partial O$: normalizing $J$ proves the first property,
and $a_h$ preserves $O$ and its boundary. Its orbit
is the same $Lz=Gz$ as that of $z$, so it has the
same intrinsic type. Thus pullback by $h$ permutes
the full set $\mathscr F$. The inverse $h^{-1}$
gives the inverse permutation. After lifting to $X$,
the two sides of (C.46.5) differ at most by the
specified target deck translation, which leaves
their fiber partitions identical.

Every lift $\widetilde h$ of $h|_O$ consequently
permutes the two coordinate foliations of $D$.
It is a coordinate pair of line homeomorphisms,
possibly followed by exchange of the coordinates.
For clarity, preservation of a fiber partition
defines its map on the corresponding real quotient;
the open quotient projection proves continuity of
that map and of its inverse. Surjectivity of both
coordinate projections then makes it a line
homeomorphism. It may be increasing or decreasing.
These quotient maps give exactly the stated
coordinate form of $\widetilde h$.

Moreover $\widetilde h$ commutes with $a_t$.
Their commutator projects to the identity on $O$
and is a continuous deck-valued function of $t$,
equal to the identity at $t=0$. The discrete deck
group makes it the identity for every $t$.
C.27's naturality under every line homeomorphism,
together with the identical construction in both
coordinates, now gives
\[
      \widetilde h\widetilde K_{s,t}\widetilde h^{-1}
                           =\widetilde K_{s,t}.         \tag{C.46.6}
\]
In particular this holds for every source deck
transformation, which is a lift of $h=\id$.
Consequently the homotopy descends to $K_{s,t}$
on $O$. The evaluation maps and their inverse
maps descend continuously through the open covering
quotient. This proves compact-open continuity,
the endpoint identities, and the asserted
commutation with the original operators.

Finally, take any other specified $z$. Its lifted
fiber partition equals one of the two chosen ones.
The first alternative of C.45 identifies its real
map as a line homeomorphism composed with that
coordinate. Naturality of the median formula under
this exact reparametrization proves
\[
 F_z\widetilde K_{s,t}=\widetilde Q^{B_z}_{s,t}F_z.
\]
For a circle target the formula commutes with its
integer deck translations, so it descends; for a
line target it already is the stated target map.
Using $q_zF_z=\Phi_zp$ proves (C.46.2).
No ambient limit of points approaching $\partial O$
has been inferred from these quotient equations.
\end{proof}
